% ============================================================================
% II. TÓM TẮT
% ============================================================================
\section{TÓM TẮT}

Các hệ thống Retrieval-Augmented Generation truyền thống thường cắt văn bản theo cửa sổ token cố định, làm đứt gãy bảng biểu và khối mã nguồn, đánh mất tiêu đề mục cha của các tiểu mục trùng tên và làm lệch số thứ tự dòng so với tệp gốc. Với sự hỗ trợ của các \texttt{Coding Agent}, người dùng có thể tạo bộ \texttt{Parser} chuyên biệt cho từng kho tài liệu nếu có một RAG framework đứng độc lập ở bên ngoài quy định hợp đồng đầu ra và cung cấp sẵn tầng lưu trữ, truy xuất lẫn đồ thị tri thức. Báo cáo này thiết kế một RAG framework phân cấp trên PostgreSQL 16 gồm: (1) \texttt{Parser interface} với 5 bất biến bắt buộc về đường dẫn phân cấp \texttt{ltree}, ánh xạ dòng 1:1, văn bản gắn tiền tố ngữ cảnh cha \texttt{contextualized\_text}, danh sách tiêu đề tổ tiên và trích dẫn nguyên văn cho cạnh đồ thị; (2) động cơ truy xuất kênh đôi hợp nhất Hybrid RRF (\texttt{Qwen3-Embedding-0.6B} rút gọn Matryoshka 512 chiều kết hợp Full-Text Search \texttt{simple}) và mô hình xếp hạng lại \texttt{mmarco-mMiniLMv2} với kênh đối sánh nguyên văn regex/trigram theo phân hoạch Venn, tự động tách tiêu đề cha ra cấp nhóm theo đúng thứ tự đọc tự nhiên và gán nhãn tin cậy qua cổng \texttt{Sparse Match Gating}; và (3) hệ quản trị Knowledge Graph cùng cơ chế xác thực trích dẫn câu trả lời qua 6 công cụ MCP. Thực nghiệm với hai bộ \texttt{Parser} mẫu trên 167 tài liệu Next.js MDX và Node.js Markdown (162.757 dòng, 5.633 \texttt{chunk}, đạt 100\% toàn vẹn dòng và khối mã) trên 100 truy vấn độc lập (\texttt{TS1-Lexical-API} và \texttt{TS2-Conceptual}) xây dựng qua 3 pha cách ly mù đạt tỷ lệ tìm thấy \texttt{Found} = 97,0\% (100,0\% trên 97/97 truy vấn có phạm vi lọc hợp lệ), \texttt{Recall@5} = 92,5\%--93,0\%, \texttt{Recall@10} = 95,0\%, \texttt{MRR@10} = 78,9\%--79,0\% và \texttt{NDCG@10} = 85,0\%, vượt trội có ý nghĩa thống kê so với truy xuất từ vựng đơn lẻ theo kiểm định McNemar ($p = 0,0156$), khẳng định hiệu quả của kiến trúc tách biệt lõi RAG với bộ \texttt{Parser} lắp ghép cho \texttt{Coding Agent}.

\vspace{0.2cm}
\noindent\textbf{Từ khóa:} Agent-Oriented RAG Framework, Parser Interface, Hierarchical Chunking, Knowledge Graph, Citation Grounding, Hybrid Search, Model Context Protocol.
